\chapter{結論と展望}
\label{chap:conclusion}

\section{本研究のまとめ}
\label{sec:summary}
本研究では，4G RANと5Gコア（5GC）の物理的な結合を解消し，世代の異なるシステムを柔軟に接続するための「S1-N2変換ゲートウェイ」を提案・実装した．
従来のモバイルネットワークは，RANとCNが密結合しており，CNを5GCへ移行するにはRANの更新・改修が前提となっていた．しかし，既存のレガシーRANを改修せずにCNのみを5GCへ先行移行する「CN先行アプローチ」を実現する標準的な手段は存在しなかった．

本研究では，以下の成果を達成した．

\begin{itemize}
	\item \textbf{疎結合アーキテクチャの確立}:
	S1AP/GTP-UとNGAP/GTP-Uのプロトコル変換を行うゲートウェイを導入することで，RANとCNのライフサイクルを分離した．これにより，RAN側の改修を一切行うことなく，5GCへの接続が可能であることを示した．

	\item \textbf{Early S1 Setup方式の考案}:
	4Gと5Gのベアラ確立手順におけるタイミングの不整合（4GはAttach時に確立，5GはRegistration後に確立）を解消するため，ゲートウェイが先行してS1ベアラを確立し，その後に5G側のPDU Sessionを確立する「Early S1 Setup」シーケンスを考案・実装した．

	\item \textbf{基本動作の確認}:
	実機（sXGP）を用いた検証により，提案方式が商用端末（UE）からの接続（位置登録，データ通信）を正しく処理できることを確認した．また，ベースライン環境との比較により，変換処理に伴う性能オーバーヘッドが軽微であることを評価した．
\end{itemize}

\section{達成目標の達成状況}
第\ref{chap:background}章で設定した3つの達成目標に対する達成状況を，設計上の達成とPoCでの実証範囲に分けて整理する．

\subsection{G1: 世代不変性に基づくプロトコルマッピングの体系化}
\label{subsec:g1}
\paragraph{設計上の達成}
第\ref{chap:proposal}章において，4G/5Gの各プロトコル仕様（S1AP/NGAP，NAS EPS/NAS 5GS，GTP-U）を比較分析し，AAA・Mobility Management・U-Planeという世代不変な機能要件に基づく変換規則を体系化した．具体的には，NAS変換規則（表\ref{tab:nas_conversion_ul}，表\ref{tab:nas_conversion_dl}），ID変換規則（表\ref{tab:id_conversion}），TEIDマッピング規則（表\ref{tab:teid_mapping}），QoSマッピング規則（表\ref{tab:qos_mapping}）として文書化した．

\paragraph{PoCでの実証}
実機検証において，Attach Request→Registration Request，Attach Accept→Registration Accept等のNAS変換，およびS1AP-ID↔NGAP-IDの変換が信号レベルで成立することを確認した．

\subsection{G2: 疎結合アーキテクチャの設計と実証}
\paragraph{設計上の達成}
第\ref{chap:proposal}章および第\ref{chap:design}章において，S1/NGインターフェースを相互変換するゲートウェイアーキテクチャを設計した．特に，4Gと5Gの接続シーケンス差異を吸収する「Early S1 Setup」方式を考案し，レガシーRANから5GCへの透過的接続を可能にする設計を確立した．

\paragraph{PoCでの実証}
sXGP基地局（eNB）とOpen5GS（5GC）を用いた実機環境において，eNBおよび5GCに一切の改修を加えることなく，UEからの位置登録およびデータ通信が正常に動作することを確認した．また，ベースライン環境との比較により，変換処理による性能オーバーヘッドが実用上無視できるレベル（RTT差約3.4 ms）であることを評価した．

\subsection{G3: 5GC固有機能への拡張性の明確化}
\label{subsec:g3}
\paragraph{設計上の達成}
第\ref{chap:proposal}章において，5GC固有機能を3層（基盤層・制御メカニズム層・応用サービス層）に分類し，層ごとに異なる対応方針（透過・マッピング・スタブ）を定義した．第\ref{chap:evaluation}章では，ネットワークスライシングやポリシー制御等の5GC固有機能への拡張に向けた設計指針と課題を明らかにした．

\paragraph{PoCでの実証}
本PoCでは基本機能の実証に注力し，5GC固有機能は実装対象外とした．ただし，単一スライス（デフォルトS-NSSAI）での動作確認により，スライシング対応への拡張が既存設計の延長線上で可能であることを示した．複数スライスの動的選択やPCF連携は今後の課題として位置づけた．

\section{研究課題に対する解決策}
第\ref{chap:introduction}章で述べた「CN先行アプローチの標準的手段の欠如」という問題に対して，本研究は以下の解決策を提示した．

\subsection{密結合モデルからの脱却}
従来の3GPP標準仕様は，RANとCNを世代ごとに同期させる「密結合モデル」を前提としており，レガシーRANを改修せずに5GCへ収容する手段が存在しなかった．

本研究では，第\ref{sec:summary}節で述べた通り，S1/NGインターフェース間のプロトコル変換を行うゲートウェイを導入することで，RANとCNのライフサイクルを分離する「疎結合アーキテクチャ」を実現し，CN先行移行という新たなパスを確立した．

\subsection{プロトコル差異の吸収}
第\ref{chap:introduction}章で指摘した「世代間で不変な基本機能」と「プロトコル仕様の差異」の問題に対して，本研究は世代共通機能（AAA，Mobility Management，U-Plane）に着目した変換規則を設計した．

第\ref{subsec:g1}項で述べた通り，実機検証によりNAS変換・ID変換・TEIDマッピング・QoSマッピングの各規則が信号レベルで成立することを確認した．この結果は，機能の共通性を利用してプロトコルの差異を吸収するというアプローチの有効性を実証するものである．

\subsection{応用可能性}
第\ref{chap:evaluation}章で論じたとおり，提案した疎結合アーキテクチャは，(1) 研究開発における低コスト5GC検証環境，(2) 商用導入におけるCN先行マイグレーション戦略，(3) NTN等の物理的制約環境への適用，といった応用が期待される．ただし，これらの応用シナリオは本実験では直接検証しておらず，設計特性から導かれる可能性として位置づけられる．

\section{今後の課題}

\subsection{スケーラビリティの向上}
第\ref{chap:evaluation}章では，C-Plane/U-Plane分離配置による独立スケーリングの可能性を考察した．C-Plane変換の負荷はセッション数に比例し，U-Plane変換の負荷はスループットに比例するという特性の違いを活かした効率的なリソース配分が期待できる．

ただし，本PoCでは分離配置の実装・検証は行っておらず，UEコンテキストの外部化に伴う複雑性・遅延の増加については未検証である．また，大規模トラフィックに対する性能最適化として，TEID変換テーブルの検索処理やパケットのコピー処理において，DPDK（Data Plane Development Kit）\cite{dpdk}やeBPF（extended Berkeley Packet Filter）\cite{ebpf}等の高速化技術を導入する余地がある．

\subsection{移動管理機能の拡張}
第\ref{chap:evaluation}章で述べた通り，本PoCでは移動管理機能のうち初期登録（Attach/Registration）の変換規則が信号レベルで成立することを確認したが，ハンドオーバー等の移動管理機能は設計のみで実証は行っていない．実運用においては基地局間の移動が必須となるケースも存在するため，S1ハンドオーバー手順とN2ハンドオーバー手順の相互変換ロジックを実装し，移動透過性を確保することが次のステップとなる．

\subsection{NTN環境での実証実験}
第\ref{chap:evaluation}章では，NTNへの適用可能性として，衛星やHAPS上のレガシーRANを地上のCNのみの更新で次世代CNへ収容できる可能性を考察した．また，提案ゲートウェイがSCTP接続を終端することで，衛星区間のタイムアウト値を独立に最適化できる特性を示した．

ただし，高遅延条件（長RTT）を付加した評価や実際の衛星回線を用いた実証実験は実施しておらず，衛星回線エミュレータやフィールド実験を通じて提案方式のロバスト性を検証・強化していく必要がある．
\subsection{5GC固有機能への対応}
第\ref{chap:evaluation}章では，5GC固有機能を3層（基盤層・制御メカニズム層・応用サービス層）に分類し，層ごとに異なる対応方針（透過・マッピング・スタブ）を設計した．PoCでは単一スライス（デフォルトS-NSSAI）での基本動作を確認したが，複数スライスの動的選択やPCF連携は検証対象外とした．

スライシングに対応するためには，S1APのQCIとNGAPのS-NSSAI/5QIを動的にマッピングする機能拡張が必要である．また，PCFとの連携によるポリシ制御やセッション単位の課金連携など，5GCの高度な機能を活用するための拡張も今後の課題である．
\subsection{6Gコアへの拡張}
第\ref{chap:evaluation}章では，本研究の変換アプローチが「将来世代（6G以降）で必要となる機能を現時点で予測し設計に織り込むことは困難である」という現実的制約を回避し，既存システムに即時適用可能な解決策を提供することを論じた．また，「世代不変機能の特定」と「プロトコル要素間の対応関係の分析」は，将来のプロトコル設計への基礎的知見としても活用できることを示した．

本研究では4G eNBと5Gコアの相互接続を実証したが，提案する疎結合アーキテクチャの汎用性を検証するため，将来的には5G RANと6Gコアの相互接続への適用可能性を検討する必要がある．認証・登録管理・QoS制御といった本質的な機能は世代を超えて継承されるという本研究の仮説が6G時代においても成立するかを検証することは，学術的にも実用的にも重要な課題である．
